\documentclass{beamer}
\usetheme{Madrid}
\usecolortheme{default}
\usepackage[utf8]{inputenc}
\usepackage{amsmath}
\usepackage{amsfonts}
\usepackage{amssymb}
\usepackage{tikz}
\usepackage{xcolor}
\usepackage{hyperref}

\title{CRISC Exam - Domain 4: Technology and Security}
\subtitle{Hard Questions with Explanations}
\author{Compiled from ISACA CRISC QAE}
\date{\today}

\begin{document}
	
	\begin{frame}
		\titlepage
	\end{frame}
	
	\begin{frame}{Outline}
		\tableofcontents
	\end{frame}
	
	% ================ SECTION 1: SECURITY CONCEPTS (CIA, IAAA, ENCRYPTION) ================
	
	\section{Security Concepts: CIA, IAAA \& Encryption}
	
	\begin{frame}{Q1: Resilient Business Processes}
		\begin{block}{Question}
			Which business requirement BEST relates to the need for resilient business and information systems processes?
		\end{block}
		\begin{enumerate}[(A)]
			\item Effectiveness
			\item Confidentiality
			\item Integrity
			\item Availability
		\end{enumerate}
	\end{frame}
	
	\begin{frame}{Q1: Explanation}
		\begin{block}{Correct Answer: D}
			Availability relates to information being available when required by the business process now and in the future. Resilience is the ability to provide and maintain an acceptable level of service during disasters or when facing operational challenges.
		\end{block}
		\begin{block}{Option Analysis}
			\begin{itemize}
				\item \textbf{A - Incorrect:} Effectiveness deals with information being relevant and delivered timely
				\item \textbf{B - Incorrect:} Confidentiality deals with protection from unauthorized disclosure
				\item \textbf{C - Incorrect:} Integrity relates to accuracy and completeness
				\item \textbf{D - Correct:} Availability is directly linked to resilience
			\end{itemize}
		\end{block}
	\end{frame}
	
	\begin{frame}{Q2: Weather Forecast System Priority}
		\begin{block}{Question}
			An information system that processes weather forecasts for public consumption is MOST likely to place its highest priority on:
		\end{block}
		\begin{enumerate}[(A)]
			\item Nonrepudiation
			\item Confidentiality
			\item Integrity
			\item Availability
		\end{enumerate}
	\end{frame}
	
	\begin{frame}{Q2: Explanation}
		\begin{block}{Correct Answer: C}
			A system that delivers weather forecasts is likely to place its highest priority on the integrity of the data. The risk practitioner should keep in mind that whether a forecast turns out to be accurate in its prediction is distinct from whether the data was accurately represented.
		\end{block}
		\begin{block}{Option Analysis}
			\begin{itemize}
				\item \textbf{A - Incorrect:} Nonrepudiation is unlikely to be important for weather forecasts
				\item \textbf{B - Incorrect:} Keeping data confidential would be at odds with public use
				\item \textbf{C - Correct:} Integrity of data is paramount
				\item \textbf{D - Incorrect:} Availability is important but integrity is higher priority
			\end{itemize}
		\end{block}
	\end{frame}
	
	\begin{frame}{Q3: Confidentiality Breach}
		\begin{block}{Question}
			Which is an example of a breach of confidentiality?
		\end{block}
		\begin{enumerate}[(A)]
			\item Data being modified by unauthorized user
			\item Data being unavailable when needed
			\item Data being disclosed to unauthorized recipient
			\item Data being destroyed
		\end{enumerate}
	\end{frame}
	
	\begin{frame}{Q3: Explanation}
		\begin{block}{Correct Answer: C}
			A breach of confidentiality means the improper disclosure of information to a recipient not authorized access by the information owner.
		\end{block}
		\begin{block}{Option Analysis}
			\begin{itemize}
				\item \textbf{A - Incorrect:} This is a breach of integrity
				\item \textbf{B - Incorrect:} This is a breach of availability
				\item \textbf{C - Correct:} Unauthorized disclosure is confidentiality breach
				\item \textbf{D - Incorrect:} This is a breach of integrity/availability
			\end{itemize}
		\end{block}
	\end{frame}
	
	\begin{frame}{Q4: CIA Triad Trade-off}
		\begin{block}{Question}
			Increasing confidentiality typically:
		\end{block}
		\begin{enumerate}[(A)]
			\item Increases availability
			\item Decreases processing time
			\item Increases processing time and may reduce availability
			\item Has no effect on availability
		\end{enumerate}
	\end{frame}
	
	\begin{frame}{Q4: Explanation}
		\begin{block}{Correct Answer: C}
			Increasing confidentiality typically increases processing time, which reduces availability. There are trade-offs between the three principles.
		\end{block}
		\begin{block}{Option Analysis}
			\begin{itemize}
				\item \textbf{A - Incorrect:} Usually decreases availability
				\item \textbf{B - Incorrect:} Usually increases processing time
				\item \textbf{C - Correct:} Trade-off between confidentiality and availability
				\item \textbf{D - Incorrect:} There is typically an effect
			\end{itemize}
		\end{block}
	\end{frame}
	
	\begin{frame}{Q5: Best Control to Prevent Unauthorized Access}
		\begin{block}{Question}
			BEST control to prevent unauthorized access to information is user:
		\end{block}
		\begin{enumerate}[(A)]
			\item Accountability
			\item Authentication
			\item Identification
			\item Access rules
		\end{enumerate}
	\end{frame}
	
	\begin{frame}{Q5: Explanation}
		\begin{block}{Correct Answer: B}
			Authentication verifies the user's identity and the right to access information according to the access rules.
		\end{block}
		\begin{block}{Option Analysis}
			\begin{itemize}
				\item \textbf{A - Incorrect:} Accountability maps activity back to responsible party
				\item \textbf{B - Correct:} Authentication verifies identity
				\item \textbf{C - Incorrect:} Identification without authentication does not grant access
				\item \textbf{D - Incorrect:} Access rules without identification and authentication do not grant access
			\end{itemize}
		\end{block}
	\end{frame}
	
	\begin{frame}{Q6: Controlling Access to Sensitive Information}
		\begin{block}{Question}
			Which control BEST protects from unauthorized individuals gaining access to sensitive information?
		\end{block}
		\begin{enumerate}[(A)]
			\item Using a challenge response system
			\item Forcing periodic password changes
			\item Monitoring unsuccessful login attempts
			\item Providing access on a need-to-know basis
		\end{enumerate}
	\end{frame}
	
	\begin{frame}{Q6: Explanation}
		\begin{block}{Correct Answer: D}
			Physical or logical system access should be assigned on a need-to-know basis (legitimate business requirements) and in ways that incorporate least privilege and segregation of duties.
		\end{block}
		\begin{block}{Option Analysis}
			\begin{itemize}
				\item \textbf{A - Incorrect:} Challenge response doesn't address access design
				\item \textbf{B - Incorrect:} Password changes don't guarantee appropriate access
				\item \textbf{C - Incorrect:} Monitoring doesn't address appropriate access rights
				\item \textbf{D - Correct:} Need-to-know basis is fundamental
			\end{itemize}
		\end{block}
	\end{frame}
	
	\begin{frame}{Q7: Best Defense Against Phishing}
		\begin{block}{Question}
			Which defense is BEST to use against phishing attacks?
		\end{block}
		\begin{enumerate}[(A)]
			\item Intrusion detection system
			\item Spam filters
			\item End-user awareness
			\item Application hardening
		\end{enumerate}
	\end{frame}
	
	\begin{frame}{Q7: Explanation}
		\begin{block}{Correct Answer: C}
			Phishing attacks are a type of social engineering attack and are best defended by end-user awareness training.
		\end{block}
		\begin{block}{Option Analysis}
			\begin{itemize}
				\item \textbf{A - Incorrect:} IDS doesn't protect against phishing
				\item \textbf{B - Incorrect:} Spam filters are not as effective as user awareness
				\item \textbf{C - Correct:} User awareness is the best defense
				\item \textbf{D - Incorrect:} Application hardening doesn't protect against phishing
			\end{itemize}
		\end{block}
	\end{frame}
	
	\begin{frame}{Q8: Symmetric Encryption}
		\begin{block}{Question}
			Symmetric encryption uses:
		\end{block}
		\begin{enumerate}[(A)]
			\item Two mathematically related keys
			\item The same secret key for encryption and decryption
			\item Public and private key pairs
			\item Digital certificates
		\end{enumerate}
	\end{frame}
	
	\begin{frame}{Q8: Explanation}
		\begin{block}{Correct Answer: B}
			Symmetric key cryptography occurs when an encryption algorithm uses the same secret key to encrypt the plaintext into ciphertext and to decrypt the ciphertext back into plaintext.
		\end{block}
		\begin{block}{Option Analysis}
			\begin{itemize}
				\item \textbf{A - Incorrect:} This is asymmetric encryption
				\item \textbf{B - Correct:} Same key for both operations
				\item \textbf{C - Incorrect:} This is asymmetric encryption
				\item \textbf{D - Incorrect:} Certificates are for identity verification
			\end{itemize}
		\end{block}
	\end{frame}
	
	\begin{frame}{Q9: Digital Signature}
		\begin{block}{Question}
			A digital signature provides:
		\end{block}
		\begin{enumerate}[(A)]
			\item Confidentiality only
			\item Integrity and nonrepudiation
			\item Availability only
			\item Authentication only
		\end{enumerate}
	\end{frame}
	
	\begin{frame}{Q9: Explanation}
		\begin{block}{Correct Answer: B}
			A digital signature combines a hash function with asymmetric encryption to verify the author's identity. It provides integrity (message not changed) and nonrepudiation (sender cannot deny sending).
		\end{block}
		\begin{block}{Option Analysis}
			\begin{itemize}
				\item \textbf{A - Incorrect:} Does not provide confidentiality on its own
				\item \textbf{B - Correct:} Provides integrity and nonrepudiation
				\item \textbf{C - Incorrect:} Not related to availability
				\item \textbf{D - Incorrect:} Provides more than authentication
			\end{itemize}
		\end{block}
	\end{frame}
	
	\begin{frame}{Q10: Asymmetric Encryption Advantage}
		\begin{block}{Question}
			Asymmetric encryption's main advantage over symmetric is:
		\end{block}
		\begin{enumerate}[(A)]
			\item Faster processing speed
			\item Easier key management for secure exchange
			\item Lower computational requirements
			\item Simpler implementation
		\end{enumerate}
	\end{frame}
	
	\begin{frame}{Q10: Explanation}
		\begin{block}{Correct Answer: B}
			Asymmetric algorithms overcome the weakness of scalability of symmetric key cryptography by allowing secure key exchange. The public key can be freely distributed, making key management easier.
		\end{block}
		\begin{block}{Option Analysis}
			\begin{itemize}
				\item \textbf{A - Incorrect:} Asymmetric is slower
				\item \textbf{B - Correct:} Easier key management
				\item \textbf{C - Incorrect:} Higher computational requirements
				\item \textbf{D - Incorrect:} More complex implementation
			\end{itemize}
		\end{block}
	\end{frame}
	
	% ================ SECTION 2: NETWORKING, ACCESS \& ARCHITECTURE ================
	
	\section{Networking, Access \& Architecture}
	
	\begin{frame}{Q11: Control to Reduce Impact of Compromised Admin Account}
		\begin{block}{Question}
			Which control can reduce potential impact of malicious hacker who gains access to administrator account?
		\end{block}
		\begin{enumerate}[(A)]
			\item Multifactor authentication
			\item Audit logging
			\item Least privilege
			\item Password policy
		\end{enumerate}
	\end{frame}
	
	\begin{frame}{Q11: Explanation}
		\begin{block}{Correct Answer: C}
			Least privilege can reduce the impact of a compromised account within the scope of its intended use and limit impact to the enterprise as a whole, without restricting the performance of employees in administrative roles.
		\end{block}
		\begin{block}{Option Analysis}
			\begin{itemize}
				\item \textbf{A - Incorrect:} Control has already been bypassed
				\item \textbf{B - Incorrect:} Audit logging is lagging indicator
				\item \textbf{C - Correct:} Least privilege limits scope of damage
				\item \textbf{D - Incorrect:} Password changes may not be frequent enough
			\end{itemize}
		\end{block}
	\end{frame}
	
	\begin{frame}{Q12: Firewall Compliance Assurance}
		\begin{block}{Question}
			Which provides BEST assurance that a firewall is configured in compliance with enterprise's security policy?
		\end{block}
		\begin{enumerate}[(A)]
			\item Reviewing actual procedures
			\item Interviewing firewall administrator
			\item Reviewing parameter settings
			\item Reviewing log file for recent attacks
		\end{enumerate}
	\end{frame}
	
	\begin{frame}{Q12: Explanation}
		\begin{block}{Correct Answer: C}
			A review of the parameter settings provides a good basis for comparison of the actual configuration to the security policy and reliable audit evidence documentation.
		\end{block}
		\begin{block}{Option Analysis}
			\begin{itemize}
				\item \textbf{A - Incorrect:} Procedures show how it should be managed, not compliance
				\item \textbf{B - Incorrect:} Interview provides overview, not confirmation
				\item \textbf{C - Correct:} Parameter settings show actual configuration
				\item \textbf{D - Incorrect:} Log review shows logging is enabled, not configuration
			\end{itemize}
		\end{block}
	\end{frame}
	
	\begin{frame}{Q13: Firewall Deviation Detection}
		\begin{block}{Question}
			Which activity should a risk professional perform to determine whether firewall deployments are deviating from security policy?
		\end{block}
		\begin{enumerate}[(A)]
			\item Review firewall parameter settings
			\item Review firewall intrusion prevention system logs
			\item Review firewall hardening procedures
			\item Analyze firewall log file for recent attacks
		\end{enumerate}
	\end{frame}
	
	\begin{frame}{Q13: Explanation}
		\begin{block}{Correct Answer: A}
			Firewall parameter settings will tie in with the configurations linked to the governing security policy. If the parameter settings differ from what policy states or requires, then there is a deviation.
		\end{block}
		\begin{block}{Option Analysis}
			\begin{itemize}
				\item \textbf{A - Correct:} Parameters show compliance with policy
				\item \textbf{B - Incorrect:} IPS logs show blocked packets, not configuration
				\item \textbf{C - Incorrect:} Hardening procedures show what was expected, not actual
				\item \textbf{D - Incorrect:} Attacks may not be covered in security policy
			\end{itemize}
		\end{block}
	\end{frame}
	
	\begin{frame}{Q14: Best Defense Against External Network Attack}
		\begin{block}{Question}
			Which is BEST way to ensure corporate network is adequately secured against external attack?
		\end{block}
		\begin{enumerate}[(A)]
			\item Use intrusion detection system
			\item Establish minimum security baselines
			\item Implement vendor recommended settings
			\item Perform periodic penetration testing
		\end{enumerate}
	\end{frame}
	
	\begin{frame}{Q14: Explanation}
		\begin{block}{Correct Answer: D}
			Penetration testing is the best way to ensure that perimeter security is adequate.
		\end{block}
		\begin{block}{Option Analysis}
			\begin{itemize}
				\item \textbf{A - Incorrect:} IDS detects but doesn't confirm security
				\item \textbf{B - Incorrect:} Baselines are beneficial but not sufficient
				\item \textbf{C - Incorrect:} Vendor settings are beneficial but not sufficient
				\item \textbf{D - Correct:} Penetration testing provides highest assurance
			\end{itemize}
		\end{block}
	\end{frame}
	
	\begin{frame}{Q15: Penetration Test Preparation}
		\begin{block}{Question}
			Which should be in place before a black box penetration test begins?
		\end{block}
		\begin{enumerate}[(A)]
			\item Clearly stated definition of scope
			\item Previous test results
			\item Proper communication and awareness training
			\item Incident response plan
		\end{enumerate}
	\end{frame}
	
	\begin{frame}{Q15: Explanation}
		\begin{block}{Correct Answer: A}
			A clearly stated definition of scope ensures a proper understanding of risk and success criteria.
		\end{block}
		\begin{block}{Option Analysis}
			\begin{itemize}
				\item \textbf{A - Correct:} Scope definition is essential
				\item \textbf{B - Incorrect:} Previous results may help define scope
				\item \textbf{C - Incorrect:} Communication/training not necessary
				\item \textbf{D - Incorrect:} IR plan not necessary
			\end{itemize}
		\end{block}
	\end{frame}
	
	\begin{frame}{Q16: Cloud Deployment for Collaborative Research}
		\begin{block}{Question}
			Which cloud computing deployment model is MOST appropriate for collaborative research program between universities?
		\end{block}
		\begin{enumerate}[(A)]
			\item Private cloud
			\item Community cloud
			\item Public cloud
			\item Hybrid cloud
		\end{enumerate}
	\end{frame}
	
	\begin{frame}{Q16: Explanation}
		\begin{block}{Correct Answer: B}
			A community cloud deployment model is appropriate because it is shared between several enterprises and supports a specific community that has a common mission or interest.
		\end{block}
		\begin{block}{Option Analysis}
			\begin{itemize}
				\item \textbf{A - Incorrect:} Private cloud is for single enterprise
				\item \textbf{B - Correct:} Community cloud is for shared mission
				\item \textbf{C - Incorrect:} Public cloud is for general public
				\item \textbf{D - Incorrect:} Hybrid is combination
			\end{itemize}
		\end{block}
	\end{frame}
	
	\begin{frame}{Q17: Highest Tolerance for Public Cloud}
		\begin{block}{Question}
			Senior management will MOST likely have highest tolerance for moving which to public cloud?
		\end{block}
		\begin{enumerate}[(A)]
			\item Credit-card processing
			\item Research and development
			\item Legacy financial system
			\item Corporate email system
		\end{enumerate}
	\end{frame}
	
	\begin{frame}{Q17: Explanation}
		\begin{block}{Correct Answer: D}
			Consideration for moving processes and information to the cloud should include the criticality, complexity and classification of the data. Of the options offered, the corporate email system has the least competitive distinction, complexity, and sensitive/highly classified information.
		\end{block}
		\begin{block}{Option Analysis}
			\begin{itemize}
				\item \textbf{A - Incorrect:} Credit-card processing is high security
				\item \textbf{B - Incorrect:} R\&D contains confidential proprietary information
				\item \textbf{C - Incorrect:} Legacy financial system is sensitive and complex
				\item \textbf{D - Correct:} Email has least competitive distinction
			\end{itemize}
		\end{block}
	\end{frame}
	
	\begin{frame}{Q18: Cloud Provider Greatest Risk}
		\begin{block}{Question}
			Which poses GREATEST risk to enterprise that recently engaged services of cloud provider?
		\end{block}
		\begin{enumerate}[(A)]
			\item Cloud provider's facility is in same vicinity as subscriber
			\item Service level agreement is ambiguous
			\item References from other customers were not obtained
			\item Auditing vendor requires dependence on third-party audit firm
		\end{enumerate}
	\end{frame}
	
	\begin{frame}{Q18: Explanation}
		\begin{block}{Correct Answer: B}
			If the service level agreement is ambiguous, it will be difficult to determine whether the provider complies.
		\end{block}
		\begin{block}{Option Analysis}
			\begin{itemize}
				\item \textbf{A - Incorrect:} No direct impact if continuity plan is adequate
				\item \textbf{B - Correct:} Ambiguous SLA makes compliance verification difficult
				\item \textbf{C - Incorrect:} References are important but can't assure delivery
				\item \textbf{D - Incorrect:} Inability to audit is less than desirable but allowed under SSAE 18
			\end{itemize}
		\end{block}
	\end{frame}
	
	\begin{frame}{Q19: Thin Client Interface Cloud Model}
		\begin{block}{Question}
			Which cloud service model provides applications accessible from client devices through a thin client interface (e.g., web browser)?
		\end{block}
		\begin{enumerate}[(A)]
			\item IaaS
			\item PaaS
			\item SaaS
			\item DaaS
		\end{enumerate}
	\end{frame}
	
	\begin{frame}{Q19: Explanation}
		\begin{block}{Correct Answer: C}
			SaaS (Software as a Service) offers the capability to use the provider's applications running on cloud infrastructure, accessible through a thin client interface such as a web browser.
		\end{block}
		\begin{block}{Option Analysis}
			\begin{itemize}
				\item \textbf{A - Incorrect:} IaaS provides processing, storage, networks
				\item \textbf{B - Incorrect:} PaaS provides a platform for deploying customer-created apps
				\item \textbf{C - Correct:} SaaS provides provider's applications via thin client
				\item \textbf{D - Incorrect:} DaaS is not a standard cloud service model in this context
			\end{itemize}
		\end{block}
	\end{frame}
	
	% ================ SECTION 3: SYSTEM DEVELOPMENT, DATA \& RESILIENCE ================
	
	\section{System Development, Data \& Resilience}
	
	\begin{frame}{Q20: First Step When Designing IS Controls}
		\begin{block}{Question}
			Which should be considered FIRST when designing information system controls?
		\end{block}
		\begin{enumerate}[(A)]
			\item Organizational strategic plan
			\item Existing IT environment
			\item Present IT budget
			\item IT strategic plan
		\end{enumerate}
	\end{frame}
	
	\begin{frame}{Q20: Explanation}
		\begin{block}{Correct Answer: A}
			Review of the enterprise's strategic plan is the first step in designing effective IS controls that would fit the enterprise's long-term plans.
		\end{block}
		\begin{block}{Option Analysis}
			\begin{itemize}
				\item \textbf{A - Correct:} Strategic plan is the starting point
				\item \textbf{B - Incorrect:} Useful but not the first task
				\item \textbf{C - Incorrect:} Budget is one component of strategic plan
				\item \textbf{D - Incorrect:} IT strategic plan exists to support enterprise's strategic plan
			\end{itemize}
		\end{block}
	\end{frame}
	
	\begin{frame}{Q21: Business Process Review}
		\begin{block}{Question}
			Which is the PRIMARY purpose of performing a business process review when responding to risk?
		\end{block}
		\begin{enumerate}[(A)]
			\item To ensure the correct solution is selected to integrate with the existing environment
			\item To identify all possible risk in the enterprise
			\item To replace the need for a risk assessment
			\item To document the organizational hierarchy
		\end{enumerate}
	\end{frame}
	
	\begin{frame}{Q21: Explanation}
		\begin{block}{Correct Answer: A}
			When responding to risk, the organization should review its business processes to ensure that the correct solution is selected to integrate and work effectively with the existing environment or determine whether the environment should be reengineered.
		\end{block}
		\begin{block}{Option Analysis}
			\begin{itemize}
				\item \textbf{A - Correct:} Ensures correct solution integration
				\item \textbf{B - Incorrect:} Risk identification is a separate process
				\item \textbf{C - Incorrect:} It complements, not replaces, risk assessment
				\item \textbf{D - Incorrect:} Documenting hierarchy is not the primary purpose
			\end{itemize}
		\end{block}
	\end{frame}
	
	\begin{frame}{Q22: Environments with Low RTO}
		\begin{block}{Question}
			Environments with low recovery time objectives (RTOs) typically have:
		\end{block}
		\begin{enumerate}[(A)]
			\item Constant transaction logs
			\item Highly available system architectures
			\item Tape backups
			\item Delayed replication
		\end{enumerate}
	\end{frame}
	
	\begin{frame}{Q22: Explanation}
		\begin{block}{Correct Answer: B}
			Environments with low RTOs typically have highly available system architectures (such as clustered servers), while those with low RPOs have highly available data architectures or constant transaction logs.
		\end{block}
		\begin{block}{Option Analysis}
			\begin{itemize}
				\item \textbf{A - Incorrect:} Constant transaction logs relate to low RPOs
				\item \textbf{B - Correct:} Low RTO needs highly available systems
				\item \textbf{C - Incorrect:} Tape backups relate to high RPOs
				\item \textbf{D - Incorrect:} Delayed replication relates to higher RPOs
			\end{itemize}
		\end{block}
	\end{frame}
	
	\begin{frame}{Q23: RTO and MTD Relationship}
		\begin{block}{Question}
			Which statement correctly describes the relationship between recovery time objective (RTO) and maximum tolerable downtime (MTD)?
		\end{block}
		\begin{enumerate}[(A)]
			\item RTO must be greater than MTD
			\item RTO must be less than MTD
			\item RTO and MTD must be equal
			\item RTO is unrelated to MTD
		\end{enumerate}
	\end{frame}
	
	\begin{frame}{Q23: Explanation}
		\begin{block}{Correct Answer: B}
			RTO must be less than MTD. If RTO is greater than or equal to MTD, the business cannot recover before failure.
		\end{block}
		\begin{block}{Option Analysis}
			\begin{itemize}
				\item \textbf{A - Incorrect:} Would mean recovery takes longer than the business can survive
				\item \textbf{B - Correct:} Recovery must complete before maximum tolerable downtime
				\item \textbf{C - Incorrect:} Equal would leave no margin
				\item \textbf{D - Incorrect:} They are directly related
			\end{itemize}
		\end{block}
	\end{frame}
	
	\begin{frame}{Q24: Data Classification Responsibility}
		\begin{block}{Question}
			Who is MOST likely responsible for data classification?
		\end{block}
		\begin{enumerate}[(A)]
			\item Data user
			\item Data owner
			\item Data custodian
			\item System administrator
		\end{enumerate}
	\end{frame}
	
	\begin{frame}{Q24: Explanation}
		\begin{block}{Correct Answer: B}
			The data owner is responsible for classifying data according to the enterprise's data classification scheme. The classification scheme then defines who is eligible to access the data and what controls are required.
		\end{block}
		\begin{block}{Option Analysis}
			\begin{itemize}
				\item \textbf{A - Incorrect:} Data user gains access based on business need
				\item \textbf{B - Correct:} Data owner classifies data
				\item \textbf{C - Incorrect:} Data custodian implements controls
				\item \textbf{D - Incorrect:} System administrator is a custodian
			\end{itemize}
		\end{block}
	\end{frame}
	
	\begin{frame}{Q25: Data Custodian Responsibility}
		\begin{block}{Question}
			Which role is responsible for implementing and maintaining the controls that protect data, such as backups and encryption?
		\end{block}
		\begin{enumerate}[(A)]
			\item Data owner
			\item Data custodian
			\item Data user
			\item Data subject
		\end{enumerate}
	\end{frame}
	
	\begin{frame}{Q25: Explanation}
		\begin{block}{Correct Answer: B}
			The data custodian implements controls and maintains data, performing tasks such as backup, encryption and security patching on behalf of the data owner.
		\end{block}
		\begin{block}{Option Analysis}
			\begin{itemize}
				\item \textbf{A - Incorrect:} Data owner decides what happens to data
				\item \textbf{B - Correct:} Custodian protects and maintains data
				\item \textbf{C - Incorrect:} User uses data for business purposes
				\item \textbf{D - Incorrect:} Data subject is the individual the data is about
			\end{itemize}
		\end{block}
	\end{frame}
	
	\begin{frame}{Q26: User Entitlement Changes}
		\begin{block}{Question}
			Who is responsible for assigning user entitlement changes?
		\end{block}
		\begin{enumerate}[(A)]
			\item IT department
			\item Data owner
			\item System administrator
			\item Data custodian
		\end{enumerate}
	\end{frame}
	
	\begin{frame}{Q26: Explanation}
		\begin{block}{Correct Answer: B}
			The data owner is responsible for assigning user entitlement changes. IT implements; the owner decides.
		\end{block}
		\begin{block}{Option Analysis}
			\begin{itemize}
				\item \textbf{A - Incorrect:} IT implements but does not decide
				\item \textbf{B - Correct:} Owner decides on entitlement changes
				\item \textbf{C - Incorrect:} Administrator implements, does not decide
				\item \textbf{D - Incorrect:} Custodian implements controls, does not decide access
			\end{itemize}
		\end{block}
	\end{frame}
	
	% ================ SECTION 4: EMERGING TECHNOLOGY, PRIVACY \& NEW RISK ================
	
	\section{Emerging Technology, Privacy \& New Risk}
	
	\begin{frame}{Q27: Vulnerability from Weak Passwords}
		\begin{block}{Question}
			Weak passwords and transmission over unprotected communication lines are examples of:
		\end{block}
		\begin{enumerate}[(A)]
			\item Vulnerabilities
			\item Threats
			\item Probabilities
			\item Impacts
		\end{enumerate}
	\end{frame}
	
	\begin{frame}{Q27: Explanation}
		\begin{block}{Correct Answer: A}
			Vulnerabilities represent characteristics of information resources that may be exploited by a threat.
		\end{block}
		\begin{block}{Option Analysis}
			\begin{itemize}
				\item \textbf{A - Correct:} Weak passwords and unprotected transmission are vulnerabilities
				\item \textbf{B - Incorrect:} Threats are circumstances with potential to cause harm
				\item \textbf{C - Incorrect:} Probabilities represent likelihood of occurrence
				\item \textbf{D - Incorrect:} Impacts represent outcome of threat exploiting vulnerability
			\end{itemize}
		\end{block}
	\end{frame}
	
	\begin{frame}{Q28: Lack of Adequate Controls}
		\begin{block}{Question}
			A lack of adequate controls represents:
		\end{block}
		\begin{enumerate}[(A)]
			\item A vulnerability
			\item An impact
			\item An asset
			\item A threat
		\end{enumerate}
	\end{frame}
	
	\begin{frame}{Q28: Explanation}
		\begin{block}{Correct Answer: A}
			Lack of adequate controls represents a vulnerability, exposing sensitive information and data to the risk of malicious damage, attack, or unauthorized access.
		\end{block}
		\begin{block}{Option Analysis}
			\begin{itemize}
				\item \textbf{A - Correct:} Vulnerability = lack of controls
				\item \textbf{B - Incorrect:} Impact is financial loss measure
				\item \textbf{C - Incorrect:} Assets have tangible/intangible value
				\item \textbf{D - Incorrect:} Threat is potential cause
			\end{itemize}
		\end{block}
	\end{frame}
	
	\begin{frame}{Q29: Threat vs. Risk Misclassification}
		\begin{block}{Question}
			Which of the following is an example of a threat that is often misclassified as a risk?
		\end{block}
		\begin{enumerate}[(A)]
			\item Data breach resulting in fines
			\item Ransomware attack
			\item Loss of customer trust
			\item Regulatory penalty
		\end{enumerate}
	\end{frame}
	
	\begin{frame}{Q29: Explanation}
		\begin{block}{Correct Answer: B}
			A ransomware attack is a threat event. It becomes a risk only when combined with the probability of it occurring and the impact it would have on the organization. A common mistake is improperly categorizing threat events as "risks."
		\end{block}
		\begin{block}{Option Analysis}
			\begin{itemize}
				\item \textbf{A - Incorrect:} This describes an impact/consequence
				\item \textbf{B - Correct:} Ransomware is a threat until it impacts something important
				\item \textbf{C - Incorrect:} Loss of trust is an impact
				\item \textbf{D - Incorrect:} Regulatory penalty is an impact
			\end{itemize}
		\end{block}
	\end{frame}
	
	\begin{frame}{Q30: Quantum Computing Impact}
		\begin{block}{Question}
			The MOST critical concern related to quantum computing for all organizations is its impact on:
		\end{block}
		\begin{enumerate}[(A)]
			\item Job displacement
			\item Applied cryptography
			\item Supply chain vulnerabilities
			\item Power generation sustainability
		\end{enumerate}
	\end{frame}
	
	\begin{frame}{Q30: Explanation}
		\begin{block}{Correct Answer: B}
			One of the most critical concerns, impactful to all organizations, is quantum computing's impact on applied cryptography. Encryption used to secure financial transactions, protect sensitive information and safeguard digital identities will effectively become obsolete once quantum computers become mainstream.
		\end{block}
		\begin{block}{Option Analysis}
			\begin{itemize}
				\item \textbf{A - Incorrect:} Job displacement is a concern but not the most universal
				\item \textbf{B - Correct:} Cryptography impact affects all organizations
				\item \textbf{C - Incorrect:} Supply chain is a concern but not the most critical
				\item \textbf{D - Incorrect:} Sustainability is a concern but not the most critical
			\end{itemize}
		\end{block}
	\end{frame}
	
	\begin{frame}{Q31: Deepfake Social Engineering}
		\begin{block}{Question}
			Which emerging technology poses a substantial challenge for organizations that authorize financial transfers based on verbal approval of an authority figure?
		\end{block}
		\begin{enumerate}[(A)]
			\item Big data analytics
			\item Deepfake technology
			\item Blockchain
			\item Internet of Things
		\end{enumerate}
	\end{frame}
	
	\begin{frame}{Q31: Explanation}
		\begin{block}{Correct Answer: B}
			Deepfake technology can create convincing audio and video imitations of a person based on samples. This poses substantial challenges for organizations that authorize financial transfers, release data or carry out other sensitive tasks based on verbal or visual approval of an authority figure.
		\end{block}
		\begin{block}{Option Analysis}
			\begin{itemize}
				\item \textbf{A - Incorrect:} Big data analytics is not a deception technology
				\item \textbf{B - Correct:} Deepfakes enable convincing impersonation
				\item \textbf{C - Incorrect:} Blockchain is a ledger technology
				\item \textbf{D - Incorrect:} IoT is about connected devices
			\end{itemize}
		\end{block}
	\end{frame}
	
	\begin{frame}{Q32: AI Risk Management}
		\begin{block}{Question}
			Which of the following is an example of a nontechnical risk that enterprises may overlook when using AI solutions?
		\end{block}
		\begin{enumerate}[(A)]
			\item Software bugs
			\item Biased or incorrect outputs impacting human rights and safety
			\item Unpatched vulnerabilities
			\item Network latency
		\end{enumerate}
	\end{frame}
	
	\begin{frame}{Q32: Explanation}
		\begin{block}{Correct Answer: B}
			AI presents numerous areas of nontechnical risk, which many enterprises may overlook—namely, ethical and environmental impacts of using AI solutions. Biased or incorrect outputs that directly impact human rights and safety can result in reputational and legal damage.
		\end{block}
		\begin{block}{Option Analysis}
			\begin{itemize}
				\item \textbf{A - Incorrect:} Software bugs are technical risk
				\item \textbf{B - Correct:} Ethical/bias impacts are nontechnical risk
				\item \textbf{C - Incorrect:} Unpatched vulnerabilities are technical risk
				\item \textbf{D - Incorrect:} Network latency is technical risk
			\end{itemize}
		\end{block}
	\end{frame}
	
	\begin{frame}{Q33: Informed Consent}
		\begin{block}{Question}
			Which privacy concept requires that data be collected, used and retained with the agreement of the individual, and that the individual be able to revoke that agreement?
		\end{block}
		\begin{enumerate}[(A)]
			\item Minimization
			\item Purpose limitation
			\item Informed consent
			\item Destruction
		\end{enumerate}
	\end{frame}
	
	\begin{frame}{Q33: Explanation}
		\begin{block}{Correct Answer: C}
			Data subject to privacy regulations should be collected, used and retained with the informed consent of the subject. There may also need to be a process for revocation of consent.
		\end{block}
		\begin{block}{Option Analysis}
			\begin{itemize}
				\item \textbf{A - Incorrect:} Minimization is about collecting only what is needed
				\item \textbf{B - Incorrect:} Purpose limitation restricts use to the stated purpose
				\item \textbf{C - Correct:} Informed consent includes revocation rights
				\item \textbf{D - Incorrect:} Destruction is about eliminating data when no longer needed
			\end{itemize}
		\end{block}
	\end{frame}
	
	\begin{frame}{Q34: Privacy Impact Assessment}
		\begin{block}{Question}
			A privacy impact assessment (PIA) primarily accounts for the effects of a compromise on the:
		\end{block}
		\begin{enumerate}[(A)]
			\item Enterprise
			\item Data subject
			\item IT department
			\item Regulator
		\end{enumerate}
	\end{frame}
	
	\begin{frame}{Q34: Explanation}
		\begin{block}{Correct Answer: B}
			Similar to a risk assessment, a PIA accounts for the effects of a compromise on the subject rather than impact to the enterprise.
		\end{block}
		\begin{block}{Option Analysis}
			\begin{itemize}
				\item \textbf{A - Incorrect:} Enterprise impact is the focus of risk assessment
				\item \textbf{B - Correct:} PIA focuses on harm to the individual
				\item \textbf{C - Incorrect:} IT department is not the focus
				\item \textbf{D - Incorrect:} Regulator is not the focus
			\end{itemize}
		\end{block}
	\end{frame}
	
	\begin{frame}{Q35: Data Minimization}
		\begin{block}{Question}
			Why is data minimization a key principle for managing privacy risk?
		\end{block}
		\begin{enumerate}[(A)]
			\item It reduces storage costs
			\item It limits the potential impact of a data breach
			\item It improves system performance
			\item It simplifies backup procedures
		\end{enumerate}
	\end{frame}
	
	\begin{frame}{Q35: Explanation}
		\begin{block}{Correct Answer: B}
			Data is both an asset and a potential disaster for enterprises subject to strong privacy regulations should a data breach occur. Privacy is best assured when only data that is expressly needed and relevant for a particular purpose is collected, limiting the impact of a breach.
		\end{block}
		\begin{block}{Option Analysis}
			\begin{itemize}
				\item \textbf{A - Incorrect:} Cost reduction is a benefit but not the key privacy rationale
				\item \textbf{B - Correct:} Less data means less exposure in a breach
				\item \textbf{C - Incorrect:} Performance is a secondary benefit
				\item \textbf{D - Incorrect:} Backup simplification is a secondary benefit
			\end{itemize}
		\end{block}
	\end{frame}
	
	\begin{frame}{Q36: What is a Risk Event?}
		\begin{block}{Question}
			For something to be a risk event, it must involve three key factors. Which of the following correctly lists them?
		\end{block}
		\begin{enumerate}[(A)]
			\item Asset, control, audit
			\item Threat event, probability, impact
			\item Vulnerability, control, impact
			\item Threat, asset, control
		\end{enumerate}
	\end{frame}
	
	\begin{frame}{Q36: Explanation}
		\begin{block}{Correct Answer: B}
			For something to be a risk event, it must involve three key factors: 1) a threat event, 2) probability (likelihood of the outcome), and 3) impact (magnitude of loss resulting from a threat exploiting a vulnerability).
		\end{block}
		\begin{block}{Option Analysis}
			\begin{itemize}
				\item \textbf{A - Incorrect:} Controls and audits are not key factors of a risk event
				\item \textbf{B - Correct:} Threat event, probability, impact
				\item \textbf{C - Incorrect:} A control is not a factor; vulnerability alone is not a risk event
				\item \textbf{D - Incorrect:} Control is not a key factor
			\end{itemize}
		\end{block}
	\end{frame}
	
	\begin{frame}{Q37: Vulnerability Assessment vs. Penetration Test}
		\begin{block}{Question}
			What is the key difference between a vulnerability assessment and a penetration test?
		\end{block}
		\begin{enumerate}[(A)]
			\item A vulnerability assessment is more expensive
			\item A penetration test actively attempts to exploit vulnerabilities to determine if a breach is possible
			\item A vulnerability assessment requires management approval
			\item A penetration test only identifies known vulnerabilities
		\end{enumerate}
	\end{frame}
	
	\begin{frame}{Q37: Explanation}
		\begin{block}{Correct Answer: B}
			A vulnerability assessment identifies potential points of compromise. A penetration test simulates an adversary using intrusive methods in an attempt to circumvent controls and exploit potential vulnerabilities to achieve objectives.
		\end{block}
		\begin{block}{Option Analysis}
			\begin{itemize}
				\item \textbf{A - Incorrect:} Cost is not the key differentiator
				\item \textbf{B - Correct:} Penetration testing actively exploits
				\item \textbf{C - Incorrect:} Both require management approval
				\item \textbf{D - Incorrect:} Vulnerability scans identify known issues; pen tests find unknown
			\end{itemize}
		\end{block}
	\end{frame}
	
	\begin{frame}{Q38: CVSS Metrics}
		\begin{block}{Question}
			Which Common Vulnerability Scoring System (CVSS) metric group can be customized based on the vulnerability's impact within a specific organizational context?
		\end{block}
		\begin{enumerate}[(A)]
			\item Base metrics
			\item Temporal metrics
			\item Environmental metrics
			\item Foundational metrics
		\end{enumerate}
	\end{frame}
	
	\begin{frame}{Q38: Explanation}
		\begin{block}{Correct Answer: C}
			Environmental metrics customize the score based on the vulnerability's impact within a specific organizational context. Base metrics define inherent characteristics, and temporal metrics reflect the current state of exploitability.
		\end{block}
		\begin{block}{Option Analysis}
			\begin{itemize}
				\item \textbf{A - Incorrect:} Base metrics are constant across environments
				\item \textbf{B - Incorrect:} Temporal metrics change over time but are not org-specific
				\item \textbf{C - Correct:} Environmental metrics are org-specific
				\item \textbf{D - Incorrect:} Foundational is not a CVSS metric group
			\end{itemize}
		\end{block}
	\end{frame}
	
	\begin{frame}{Q39: Security Information and Event Management (SIEM)}
		\begin{block}{Question}
			What is the primary function of a SIEM system?
		\end{block}
		\begin{enumerate}[(A)]
			\item To store backup data
			\item To correlate data from multiple sources to identify security events
			\item To manage user identities
			\item To encrypt network traffic
		\end{enumerate}
	\end{frame}
	
	\begin{frame}{Q39: Explanation}
		\begin{block}{Correct Answer: B}
			SIEM systems are integrated data correlation tools that capture data from multiple sources and analyze system, application and network activity for possible security events.
		\end{block}
		\begin{block}{Option Analysis}
			\begin{itemize}
				\item \textbf{A - Incorrect:} Not for backup storage
				\item \textbf{B - Correct:} Correlation of multiple data sources
				\item \textbf{C - Incorrect:} Identity management is separate
				\item \textbf{D - Incorrect:} Not for traffic encryption
			\end{itemize}
		\end{block}
	\end{frame}
	
	\begin{frame}{Q40: Endpoint Detection and Response (EDR)}
		\begin{block}{Question}
			Unlike signature-based antivirus, Endpoint Detection and Response (EDR) is primarily designed to:
		\end{block}
		\begin{enumerate}[(A)]
			\item Detect known attack patterns
			\item Be predictive by identifying threats designed to evade traditional defenses
			\item Block all malware
			\item Store endpoint logs indefinitely
		\end{enumerate}
	\end{frame}
	
	\begin{frame}{Q40: Explanation}
		\begin{block}{Correct Answer: B}
			EDR attempts to be predictive in nature, focusing on identifying threats and malware designed to evade traditional security defenses. It leverages cyberthreat intelligence and ML in conjunction with file analysis.
		\end{block}
		\begin{block}{Option Analysis}
			\begin{itemize}
				\item \textbf{A - Incorrect:} Signature-based AV detects known patterns
				\item \textbf{B - Correct:} EDR is predictive and finds evasive threats
				\item \textbf{C - Incorrect:} No control blocks all malware
				\item \textbf{D - Incorrect:} Log storage is a feature but not the primary purpose
			\end{itemize}
		\end{block}
	\end{frame}
	
	\begin{frame}{Q41: Technology Refresh Risk}
		\begin{block}{Question}
			Which is a risk consideration related to retaining legacy systems that are past their supported life?
		\end{block}
		\begin{enumerate}[(A)]
			\item Ability to patch/mitigate vulnerabilities may be limited
			\item Lack of documentation
			\item Vendor lock-in
			\item All of the above
		\end{enumerate}
	\end{frame}
	
	\begin{frame}{Q41: Explanation}
		\begin{block}{Correct Answer: D}
			Risk considerations related to technology refresh include age of equipment, limited ability to patch/mitigate vulnerabilities, lack of documentation, limited expertise for maintenance, vendor lock-in, and availability of replacement parts.
		\end{block}
		\begin{block}{Option Analysis}
			\begin{itemize}
				\item \textbf{A - Incorrect:} A valid concern but not the only one
				\item \textbf{B - Incorrect:} A valid concern but not the only one
				\item \textbf{C - Incorrect:} A valid concern but not the only one
				\item \textbf{D - Correct:} All of these are valid legacy system risks
			\end{itemize}
		\end{block}
	\end{frame}
	
	\begin{frame}{Q42: Zero Trust Architecture}
		\begin{block}{Question}
			A zero trust architecture focuses on protecting resources and controlling access based on:
		\end{block}
		\begin{enumerate}[(A)]
			\item The network segment the resource is on
			\item The user's identity, device posture and context
			\item The physical location of the data center
			\item The firewall rules at the perimeter
		\end{enumerate}
	\end{frame}
	
	\begin{frame}{Q42: Explanation}
		\begin{block}{Correct Answer: B}
			Zero trust moves defenses from static, network-based perimeters to focus on users, assets and resources. It requires real-time monitoring and granular visibility of user and application behavior attributes to make access decisions.
		\end{block}
		\begin{block}{Option Analysis}
			\begin{itemize}
				\item \textbf{A - Incorrect:} Zero trust does not rely on network segments
				\item \textbf{B - Correct:} Access is based on identity, device and context
				\item \textbf{C - Incorrect:} Physical location is not the primary control factor
				\item \textbf{D - Incorrect:} Perimeter firewalls are not the basis of zero trust
			\end{itemize}
		\end{block}
	\end{frame}
	
	% ================ FINAL SLIDES ================
	
	\begin{frame}{Domain 4 Key Takeaways}
		\begin{block}{Essential Concepts}
			\begin{itemize}
				\item CIA triad involves trade-offs; you can't maximize all three
				\item Authentication verifies identity; authorization determines permissions
				\item Digital signatures provide integrity and nonrepudiation
				\item Data owner classifies data; data custodian implements controls
				\item Penetration testing actively exploits vulnerabilities
				\item RTO must be less than MTD
				\item Zero trust focuses on users, assets and context, not network perimeters
				\item Emerging tech (AI, quantum, deepfakes) introduces new risk categories
			\end{itemize}
		\end{block}
	\end{frame}
	
	\begin{frame}{Final Advice}
		\begin{block}{Exam Preparation}
			\centering
			\vspace{0.3cm}
			\textbf{Remember:}
			\begin{itemize}
				\item Understand the \textit{why} behind each concept
				\item Know the difference between similar terms (vulnerability vs. threat, IDS vs. IPS)
				\item Understand the risk management life cycle (Identify, Assess, Respond, Monitor)
				\item Practice scenario-based questions
				\item Use the ISACA mindset - choose the MOST appropriate answer
			\end{itemize}
			\vspace{0.5cm}
			\textbf{Good Luck on Your CRISC Exam!}
		\end{block}
	\end{frame}
	
\end{document}